\BeleskeID{02-s047}
% element: 02-s047:e01
% element: 02-s047:e02
% element: 02-s047:e03
% element: 02-s047:d01

Prvo dvoulazno I kolo formira \(\bar D\bar B\), drugo formira \(D\bar C\), a treće množi rezultat drugog sa \(\bar A\). Zaseban čip sa troulaznim I kolom zato nije potreban. U četvorostrukom kućištu dvoulaznih I kola ostaje jedno slobodno kolo.

Putanja od \(D\) ili \(\bar C\) kroz drugi proizvod sada prolazi kroz dva I kola pre završnog ILI kola. Ako su kašnjenja uporediva, ona može biti sporija od putanje kroz jedno troulazno I i završno ILI. Asocijativnost omogućava ovu preraspodelu bez promene stabilne funkcije, ali ne čuva automatski vreme odziva. Broj kućišta je manji iako je broj upotrebljenih I kola porastao.
